%!TEX root = ../main.tex

\section{Task A: Blink an LED}

\subsection{a) Calculations:}

Gegeben sind \(T = 3\,\mathrm{s}\) und \(D = 66.6\%\).

\[
	t_{\mathrm{on}} = T \cdot D = 3\,\mathrm{s} \cdot 0.666 \approx 2\,\mathrm{s}
\]
\[
	t_{\mathrm{off}} = T - t_{\mathrm{on}} = 3\,\mathrm{s} - 2\,\mathrm{s} = 1\,\mathrm{s}
\]

\subsection{b) Implementation:}

Die LED wird über \texttt{PJ7 (SEGDP)} gesteuert. In der Endlosschleife wird der Pin gesetzt (LED an), anschließend für die Einschaltzeit gewartet und danach wieder gelöscht (LED aus). Die Zeiten werden mit \texttt{HAL\_Delay()} realisiert.

\begin{lstlisting}
  /* USER CODE BEGIN 2 */
  const uint32_t onTime = 2000;
  const uint32_t offTime = 1000;
  /* USER CODE END 2 */

  while (1)
  {
    GPIOJ->ODR |= SEGDP_Pin;
    HAL_Delay(onTime);
    GPIOJ->ODR &= ~SEGDP_Pin;
    HAL_Delay(offTime);
  }
\end{lstlisting}

\subsection{c) Results:}

Die LED blinkt mit der gewünschten Periode von \(3\,\mathrm{s}\) und ist dabei länger eingeschaltet als ausgeschaltet (ca. \(2\,\mathrm{s}\) EIN, \(1\,\mathrm{s}\) AUS).
\begin{figure}[H]
	\begin{center}
		\includegraphics[width=0.30\textwidth]{Task1.jpg}
	\end{center}
	\caption{Bord with LED dot of Segment Display on}
\end{figure}

\subsection{d) Discussion:}

Die Lösung ist einfach und übersichtlich durch feste Zeitkonstanten. \texttt{HAL\_Delay()} ist jedoch blockierend. Wiederverwendbar sind die GPIO-Steuerung und die Schleifenstruktur.
